\documentclass[11pt]{article}
\usepackage[margin=0.82in]{geometry}
\usepackage{graphicx}
\usepackage{booktabs}
\usepackage{longtable}
\usepackage{array}
\usepackage{float}
\usepackage{caption}
\usepackage{hyperref}
\usepackage{xcolor}
\usepackage{amsmath}
\usepackage{enumitem}
\usepackage{fancyvrb}
\usepackage[T1]{fontenc}
\usepackage[utf8]{inputenc}
\hypersetup{colorlinks=true, linkcolor=blue, urlcolor=blue}
\setlist[itemize]{leftmargin=1.5em}
\setlist[enumerate]{leftmargin=1.5em}
\captionsetup{font=small, labelfont=bf}
\sloppy
\title{Phase-II Robust Active-Learning Convergence, Numerical Reliability, and Optimization Report}
\date{2026-06-20}
\begin{document}
\maketitle
\tableofcontents
\newpage


\section{Executive Summary}

\begin{table}[H]
\centering
\small
\begin{tabular}{@{}p{0.43\linewidth}p{0.47\linewidth}@{}}
\toprule
Item & Value \\
\midrule
main\_phase\_converged & True \\
publication\_boundary\_audit & pass \\
publication\_ready\_for\_main\_phase\_map & True \\
frozen dataset & dataset\_iter035 \\
dataset total & 7434 \\
normal / uniform-SC / FFLO & 1867 / 715 / 4852 \\
trusted surprise & 0/127 \\
hard-risk surprise & 75/129 \\
targeted rerun count & 0 \\
\bottomrule
\end{tabular}
\caption{Final Phase-II status.}
\end{table}

Phase-II should be read as one continuous workflow: exact seed labels, surrogate training, full acquisition, robust exact-oracle labeling, trusted-label append, StopController review, and hard-risk publication audit.  The machine-learning model schedules exact BdG calls; it does not replace the thermodynamic phase criterion.

\section{From Random Seed to Closed-Loop Active Learning}

\begin{figure}[H]
\centering
\includegraphics[width=0.84\linewidth]{figures/phase2_from_seed_loop.png}
\caption{RankCapK3 active-learning loop from random seed.}
\end{figure}


The final loop contains an initial exact seed set followed by acquisition-selected batches.  Every batch is merged, filtered through trusted-label criteria, appended to the dataset, and checked against main phase-map stability and boundary-shift criteria.

\section{Optimization Timeline}

\begin{figure}[H]
\centering
\includegraphics[width=0.96\linewidth]{figures/phase2_optimization_timeline.png}
\caption{Phase-II optimization and validation timeline.}
\end{figure}


\begin{table}[H]
\centering
\small
\begin{tabular}{@{}p{0.18\linewidth}p{0.30\linewidth}p{0.42\linewidth}@{}}
\toprule
stage & objective & main\_result \\
\midrule
Initial exact grid & Reference phase structure & Coarse exact map available; active learning need is boundary refinement \\
Acquisition development & Avoid deep-normal oversampling while retaining discovery & Full acquisition retained for FFLO/uniform-SC coverage; simple-phase not used for final discovery \\
Robust exact oracle & Separate phase labels from numerical reliability & Trusted, training-eligible, rerun-required, q/Delta unresolved layers recorded \\
Incremental q-window & Reduce repeated q-window work & Incremental strip computation retained with fallback and reliability audit \\
Local-refinement audit & Find why optimized variants timed out & Mandatory overflow produced approximately 85 boxes in failed variants \\
RankCapK3 & Bound selected refinement targets & Fixed-point acceptance passed; selected targets <= 3 with physics-equivalent labels \\
One- and five-iteration validation & Check real AL behavior before full run & RankCapK3 remained stable and reduced local boxes in active-loop context \\
Full loop & Run production discovery with optimized oracle & 30 acquisition batches completed; main phase nearly stable but all-selected surprise remained high \\
Trusted-surprise tail & Use trusted labels as formal main-phase surprise gate & Stopped at iter034; trusted surprise 0/127; dataset\_iter035 frozen \\
Hard-risk publication audit & Test whether hard-risk frontier can move main boundaries & Decision A; targeted rerun count 0; publication boundary audit pass \\
\bottomrule
\end{tabular}
\caption{Stage-level objective and outcome.}
\end{table}

\section{Acquisition Function Development}

Acquisition development controlled three competing needs: discover new FFLO and uniform-SC regions, keep attention near main thermodynamic boundaries, and avoid wasting batches in deep normal interiors.  The final full profile retained boundary terms, uncertainty terms, stochastic sampling, observation repulsion, and batch repulsion.

\begin{figure}[H]
\centering
\includegraphics[width=0.84\linewidth]{figures/phase2_acquisition_oracle_flow.png}
\caption{Responsibility split between acquisition and exact oracle.}
\end{figure}


\section{Exact-Oracle Reliability Evolution}

The robust oracle records phase labels separately from reliability metadata: trusted exact status, training eligibility, rerun-required status, q-unresolved status, and Delta-unresolved status.  This separation lets the report distinguish main phase convergence from the still-active numerical frontier.

\section{Incremental q-Window and Local-Refinement Optimization}

Incremental q-window computation reduces repeated q-strip work without changing the thermodynamic phase criterion.  Local refinement was the larger bottleneck: failed cluster-optional variants overflowed to about 85 local boxes because mandatory risk targets bypassed the intended cap.  RankCapK3 fixed this by using basin-level risk annotation and a hard total cap of three refined basins.

\begin{table}[H]
\centering
\small
\begin{tabular}{@{}p{0.55\linewidth}p{0.35\linewidth}@{}}
\toprule
metric & value \\
\midrule
baseline\_local\_boxes & 6.0 \\
rankcap\_mean\_local\_boxes & 2.79296875 \\
rankcap\_max\_local\_boxes & 3 \\
baseline\_local\_refinement\_runtime\_sec & 189.767 \\
rankcap\_local\_refinement\_runtime\_sec & 88.28558737352148 \\
local\_refinement\_runtime\_reduction\_percent & 53.47684930808756 \\
baseline\_point\_total\_runtime\_sec & 234.194 \\
rankcap\_point\_total\_runtime\_sec & 117.28463183327355 \\
point\_total\_runtime\_reduction\_percent & 49.91988187858205 \\
total\_wall\_hours & 36.571320979444444 \\
wall\_minutes\_per\_exact\_iteration & 70.78320189569892 \\
\bottomrule
\end{tabular}
\caption{Optimization performance summary.}
\end{table}

\begin{figure}[H]
\centering
\includegraphics[width=0.84\linewidth]{figures/local_box_distribution.png}
\caption{Local-box distribution under RankCapK3.}
\end{figure}


\begin{figure}[H]
\centering
\includegraphics[width=0.84\linewidth]{figures/local_box_gate.png}
\caption{Local-box gate and cap behavior.}
\end{figure}


\section{Validation Ladder}

The optimized oracle was not promoted directly into a full run.  It passed fixed-point acceptance, one-iteration active-learning validation, five-iteration validation, a full-loop run, tail continuation, and the publication hard-risk boundary-impact audit.

\section{Full-Loop and Tail-Continuation Results}

\begin{table}[H]
\centering
\small
\begin{tabular}{@{}p{0.45\linewidth}p{0.30\linewidth}@{}}
\toprule
phase & count \\
\midrule
 &  \\
\bottomrule
\end{tabular}
\caption{Frozen dataset\_iter035 phase counts.}
\end{table}

\begin{table}[H]
\centering
\small
\begin{tabular}{@{}p{0.43\linewidth}p{0.47\linewidth}@{}}
\toprule
Item & Value \\
\midrule
phase\_map\_change & 0.001628727653468802 \\
normal/SC boundary shift & 0.004166666666666666 \\
uniform/FFLO boundary shift & 0.0 \\
boundary\_coverage\_p95 & 0.0046875000000001586 \\
trusted surprise & 0/127 \\
all-selected surprise & 75/256 \\
hard-risk surprise & 75/129 \\
\bottomrule
\end{tabular}
\caption{Final StopController metrics.}
\end{table}

\begin{figure}[H]
\centering
\includegraphics[width=0.84\linewidth]{figures/dataset_growth_and_phase_counts.png}
\caption{Dataset growth and phase counts.}
\end{figure}


\begin{figure}[H]
\centering
\includegraphics[width=0.84\linewidth]{figures/surrogate_metric_curves.png}
\caption{Surrogate metric curves.}
\end{figure}


\begin{figure}[H]
\centering
\includegraphics[width=0.84\linewidth]{figures/runtime_curve.png}
\caption{Runtime learning curve.}
\end{figure}


\section{Final Phase Diagram}

\begin{figure}[H]
\centering
\includegraphics[width=0.84\linewidth]{figures/final_phase_map_clean.png}
\caption{Clean final main phase map with extracted boundaries.}
\end{figure}


\begin{figure}[H]
\centering
\includegraphics[width=0.84\linewidth]{figures/final_phase_map_with_uncertainty.png}
\caption{Final phase map with hard-risk uncertainty markers.}
\end{figure}


\begin{figure}[H]
\centering
\includegraphics[width=0.84\linewidth]{figures/final_dataset_coverage.png}
\caption{Final exact dataset coverage.}
\end{figure}


\begin{figure}[H]
\centering
\includegraphics[width=0.84\linewidth]{figures/hard_risk_uncertainty_band.png}
\caption{Hard-risk uncertainty band.}
\end{figure}


\section{Convergence Logic}

Phase-map change and boundary-shift metrics measure stability of dense monitor-grid predictions between iterations.  Boundary coverage measures how densely exact samples cover the current main boundaries.  All-selected surprise is an acquisition-difficulty diagnostic, whereas trusted surprise is the clean exact-label gate for main phase convergence.

\begin{figure}[H]
\centering
\includegraphics[width=0.84\linewidth]{figures/tail_surprise_layers.png}
\caption{Trusted and hard-risk surprise layers at the tail endpoint.}
\end{figure}


\section{Hard-Risk Boundary-Impact Audit}

\begin{table}[H]
\centering
\small
\begin{tabular}{@{}p{0.65\linewidth}p{0.25\linewidth}@{}}
\toprule
reason & count \\
\midrule
rerun\_required & 110 \\
untrusted\_exact & 129 \\
training\_ineligible & 110 \\
q\_unresolved & 0 \\
delta\_unresolved & 0 \\
hard\_risk\_total & 129 \\
multi\_reason\_overlap & 110 \\
rerun\_required\_only & 0 \\
non\_rerun\_hard\_risk & 19 \\
\bottomrule
\end{tabular}
\caption{Hard-risk reason counts.}
\end{table}

\begin{table}[H]
\centering
\small
\begin{tabular}{@{}p{0.31\linewidth}p{0.14\linewidth}p{0.16\linewidth}p{0.31\linewidth}@{}}
\toprule
gate & status & value & threshold \\
\midrule
p95 <= tol & pass & 0 & 0.00416667 \\
sig. Hausdorff <= tol & pass & 0 & 0.00416667 \\
strict Hausdorff diag. & diagnostic & 0.8125 & diagnostic; sig. arc >= 0.05 \\
topology stable & pass & False & False \\
no moving cluster & pass & 0 & 0 \\
counts reconciled & pass & 129 & 129 \\
metadata present & pass & False & False \\
targeted list empty & pass & 0 & 0 \\
publication audit & pass & True & True \\
\bottomrule
\end{tabular}
\caption{Publication audit gate summary.}
\end{table}

The strict Hausdorff diagnostic reaches 0.8125 for isolated dense-grid uncertainty-marker fragments, but the significant local Hausdorff gate value is zero.  No meaningful main-boundary topology change is detected, and the targeted rerun list is empty.

\section{Supporting Report Map}

\begin{table}[H]
\centering
\small
\begin{tabular}{@{}p{0.48\linewidth}p{0.36\linewidth}@{}}
\toprule
directory & role \\
\midrule
\_supporting\_reports & supporting container \\
report\_active\_learning\_r0015\_note & main report \\
report\_phase2\_robust\_al\_final\_202606 & main report \\
acquisition\_profile\_smoke & Phase-II supporting report \\
hard\_risk\_boundary\_impact\_audit & Phase-II supporting report \\
hard\_risk\_boundary\_impact\_audit\_v2 & Phase-II supporting report \\
local\_refinement\_rankcap\_acceptance & Phase-II supporting report \\
local\_refinement\_refactor & Phase-II supporting report \\
local\_refinement\_refactor\_goal\_run & Phase-II supporting report \\
local\_refinement\_target\_construction\_dryrun & Phase-II supporting report \\
local\_refinement\_target\_logic\_audit & Phase-II supporting report \\
phase\_boundary\_stability\_audit & Phase-II supporting report \\
rankcap\_k3\_5iter\_validation & Phase-II supporting report \\
rankcap\_k3\_5iter\_validation\_recheck & Phase-II supporting report \\
rankcap\_k3\_full\_loop & Phase-II supporting report \\
rankcap\_k3\_one\_iter\_validation & Phase-II supporting report \\
\bottomrule
\end{tabular}
\caption{Main report and supporting-report map.}
\end{table}

\section{Consistency and Reproduction}

\begin{table}[H]
\centering
\small
\begin{tabular}{@{}p{0.28\linewidth}p{0.20\linewidth}p{0.24\linewidth}p{0.12\linewidth}@{}}
\toprule
check & expected & actual & status \\
\midrule
dataset\_total & 7434 & 7434 & pass \\
normal\_count & 1867 & 1867 & pass \\
uniform\_SC\_count & 715 & 715 & pass \\
FFLO\_count & 4852 & 4852 & pass \\
final\_iteration & 34 & 34 & pass \\
phase\_map\_change & <0.002 & 0.001628727653468802 & pass \\
hard\_risk\_total & 129 & 129 & pass \\
rerun\_required\_count & 110 & 110 & pass \\
local\_box\_mean & approx 2.79 & 2.79296875 & pass \\
local\_box\_max & 3 & 3 & pass \\
figure\_sources\_exist & all copied/generated & 11 & pass \\
\bottomrule
\end{tabular}
\caption{Report consistency checks.}
\end{table}

The reproduction manifest records source run IDs, dataset paths and hashes, configuration, tolerance values, input report paths, and output hashes.

\begin{table}[H]
\centering
\small
\begin{tabular}{@{}p{0.43\linewidth}p{0.47\linewidth}@{}}
\toprule
Item & Value \\
\midrule
git commit & f3c277bb419ec4108b40ebba56635ef72c2ad895 \\
candidate grid hash & f7af1e7dab418e71... \\
phase map hash & 7d96f6593d368e76... \\
\bottomrule
\end{tabular}
\caption{Reproduction identifiers.}
\end{table}

\section{Do-Not-Claim List}

\begin{enumerate}\item Do not claim the hard-risk numerical frontier has disappeared.\item Do not treat provisional hard-risk labels as definitive phase labels.\item Do not claim eta anomalies have all become physical effects.\item Do not treat topology reference curves as pointwise topology labels.\item Do not present a single-seed run as a complete multi-seed benchmark.\item Do not claim hidden-ground-truth benchmarking is complete.\item Do not treat synchronized global stress-test islands as physical boundary instabilities.\end{enumerate}

\section{Final Decision and Next Stage}

Freeze \texttt{dataset\_iter035}, RankCapK3 production oracle settings, and the trusted-surprise StopController configuration for the Phase-II main phase-map result.  Do not launch another full active-learning loop for this result.  The next physics stages are branch-resolved topology classification, hidden-ground-truth evaluation, multi-seed benchmarking, and final publication figure polishing.


\end{document}
